\documentclass[11pt,a4paper]{article}

\usepackage[margin=1in]{geometry}
\usepackage{amsmath, amssymb, amsthm}
\usepackage{mathtools}
\usepackage{bm}
\usepackage{enumitem}
\usepackage{microtype}
\usepackage[T1]{fontenc}
\usepackage{lmodern}
\usepackage[dvipsnames]{xcolor}
\usepackage{hyperref}

\hypersetup{
  colorlinks=true,
  linkcolor=blue!50!black,
  urlcolor=blue!50!black,
  citecolor=blue!50!black
}

\newtheorem{theorem}{Theorem}[section]
\newtheorem{lemma}[theorem]{Lemma}
\newtheorem{definition}[theorem]{Definition}

\title{U5: The Forced Structure of Translation Symmetry and the Momentum Operator}

\author{Allen Proxmire}

\date{April 2026}

\begin{document}

\maketitle

\begin{abstract}
We present the structural derivation of the spatial translation generator and the momentum operator on the participation-measure Hilbert space from the Event-Density (ED) primitive stack and the now-\textsc{forced} upstream items U1, U2-Discrete, U2-Continuum, and Theorem~10 (Born). Operating strictly outside Framing~1 --- i.e., \textbf{without invoking U3} in any role, derivation input or otherwise --- we decompose U5 into three sub-features: existence of a strongly continuous one-parameter unitary representation of spatial translations on the U2 Hilbert space (F1); uniqueness of its self-adjoint generator (F2); and identification of that generator with the standard momentum operator $\hat{p} = -i\hbar\nabla$, with plane waves as eigenfunctions and the standard $L^2$ Fourier transform as the unique unitary intertwiner of position and momentum representations (F3). F1 follows from Primitive~03 (graph homogeneity) and Primitive~06 (admissibility of translations along the spatial axis), with unitarity inherited from the U2 inner product. F2 follows from Stone's theorem on one-parameter unitary groups. F3 follows from the position-representation form of the translation generator together with the standard Fourier theory on $L^2(\mathbb{R}^d)$, with the principal alternatives --- fractional-Fourier, wavelet, and Mellin transforms --- explicitly dismissed by structural inconsistency with translation-driven conjugacy. All four falsifiers are dispatched. The U5 verdict is \textbf{\textsc{forced}}, with no new structural \textsc{candidate}s introduced and no dependence whatever on the sibling \textsc{candidate} U3. The arc closes the kinematic backbone of the QM-emergence Phase-1 program: U5 supplies the momentum operator that U4's Galilean integration acts upon, and the Fourier-conjugacy that Phase-1 Step~5 needs to derive the Heisenberg bound $\Delta x \cdot \Delta p \geq \hbar/2$. The U5 derivation is what makes U4's ``form-\textsc{forced}--value-\textsc{inherited}'' closure possible while preserving acyclicity. The active upstream-\textsc{candidate} inventory is unchanged by U5 alone (U5 was promoted in the 2026-04-26 cycle); the present paper documents and formalizes the result for the publication-series record.
\end{abstract}

\section{Introduction}

\subsection{What U5 is}

U5 is the structural commitment that the participation-measure Hilbert space $\mathcal{H}$ --- established by the U2-Discrete and U2-Continuum theorems --- carries a strongly continuous one-parameter unitary representation of spatial translations whose unique self-adjoint generator is the momentum operator $\hat{p}$. Equivalently, U5 supplies the kinematic translation-conjugate of position and identifies it as the operator whose eigenfunctions are plane waves and whose position-representation form is $\hat{p} = -i\hbar\nabla$.

The earlier framing of U5 in the QM-emergence Phase-1 inventory was as the \emph{adjacency-band Fourier-conjugate partition}, $b^{\mathrm{adj}} = b_x + b_p$. The two framings are equivalent: the partition exists, is two-component, is Fourier-conjugate, and is orthogonal exactly when the translation group on $\mathcal{H}$ admits a unique self-adjoint generator and the standard $L^2$ Fourier transform supplies the basis-change between its position and momentum eigenbases. The present paper adopts the translation-symmetry framing because it makes the U5 result structurally transparent: U5 is, at root, \emph{Stone's theorem applied to translation symmetry on the U2 Hilbert space}.

\subsection{Why U5 matters}

Three of the four foundational postulates of non-relativistic single-particle quantum mechanics --- the Born rule, the Bell--Tsirelson bound, and the Heisenberg uncertainty inequality --- depend on U5 either directly or via the inner-product structure that U5's translation operators preserve. The fourth foundational postulate --- Schr\"odinger evolution --- depends on U5 indirectly: the Hamiltonian $\hat{H} = \hbar^2|\hat{p}|^2/(2m) + V(\hat{x})$ derived in U4 contains the momentum operator that U5 supplies.

U5 is therefore the \emph{kinematic backbone} of the structural-foundations program. U2 builds the Hilbert space; Theorem~10 (Born) installs the probabilistic interpretation; U5 supplies the translation generator and the position-momentum conjugacy; U1 derives the participation-measure that lives in the space; U4 derives the specific Hamiltonian form that acts on it. U3, the sole remaining upstream \textsc{candidate}, supplies the time-evolution generator's existence + linear-first-order structure; U5 is its spatial-translation analog and methodologically prefigures it.

\subsection{Relationship to U1, U2, U4, and U3}

\begin{itemize}
  \item \textbf{U1.} U5 takes the participation-measure construction $P_K = \sqrt{b_K} \cdot e^{i\pi_K}$ as input. The polarity-phase factor $e^{i\pi_K}$ supplies the $U(1)$-valued phase whose gradient identifies the momentum-conjugate direction.
  \item \textbf{U2.} U5 takes the U2 Hilbert space and its sesquilinear inner product as input. Translation operators are unitary on the U2 Hilbert space; Stone's theorem applies.
  \item \textbf{U4.} U5 supplies the momentum operator $\hat{p}$ that appears in U4's Hamiltonian form. U4's Galilean Lie algebra integration acts on $\hat{p}$ and would not even be statable without it.
  \item \textbf{U3.} U5 is \emph{strictly independent} of U3. The U5 derivation invokes no time-evolution operator, no Hamiltonian, no Schr\"odinger equation. This is the central methodological feature of the present derivation and the principal sharpening over the earlier 2026-04-24 tightening-program memo, which had derived the Fourier conjugacy via U3.
\end{itemize}

\subsection{Summary of the U5 verdict}

The U5 verdict is \textbf{\textsc{forced}}. The three sub-features F1, F2, F3 close uniquely under the available structural inputs. All four falsifiers (Fal-1 through Fal-4) dispatch. No new structural \textsc{candidate} is introduced. The derivation is U3-independent. The U5 verdict is unconditional in both the discrete and continuum regimes; the U2-Continuum conformal gauge acts trivially on the translation representation and preserves the partition structure pointwise.

\subsection{Organization}

Section~\ref{sec:question} states the U5 question precisely and inventories its sub-features and falsifiers. Section~\ref{sec:inputs} inventories the structural inputs and the inputs explicitly forbidden under the U3-independence discipline. Section~\ref{sec:decomposition} decomposes U5 into F1, F2, F3 with status assignments. Section~\ref{sec:falsifiers} introduces the four falsifiers. Sections~\ref{sec:F1}--\ref{sec:F3} derive F1, F2, F3 in turn. Section~\ref{sec:falsifier-table} consolidates the falsifier dispatches. Section~\ref{sec:conditionality} documents the conditionality ledger. Section~\ref{sec:final-theorem} states the final U5 theorem. Section~\ref{sec:downstream} documents the downstream cascade. Section~\ref{sec:discussion} discusses the result's significance and identifies the U3 arc as the natural successor.

\section{Statement of the U5 Question}\label{sec:question}

\subsection{The target structure}

The target structural content of U5 is the existence and uniqueness of the self-adjoint generator of spatial translations on the U2 Hilbert space, together with its identification as the momentum operator with the standard plane-wave eigenfunctions.

Formally: let $\mathcal{H}$ denote the U2 Hilbert space of participation measures (discrete or continuum regime, gauge-class equivalent under the U2-Continuum conformal redundancy). Let $T_a$ for $a \in \mathbb{R}^d$ denote the spatial translation operator $T_a P(x) = P(x - a)$ acting componentwise on the participation-measure components. The U5 question is whether $\{T_a : a \in \mathbb{R}^d\}$ is a strongly continuous one-parameter unitary group whose unique self-adjoint generator is the standard momentum operator $\hat{p} = -i\hbar\nabla$ in the position representation, with plane waves $\langle x | k \rangle = (2\pi\hbar)^{-d/2} e^{ikx/\hbar}$ as eigenfunctions.

\subsection{The three sub-features}

\begin{itemize}
  \item \textbf{F1 (Existence of translation representation).} $\{T_a : a \in \mathbb{R}^d\}$ is a well-defined, strongly continuous, one-parameter (per spatial axis) unitary group on $\mathcal{H}$.
  \item \textbf{F2 (Uniqueness of the self-adjoint generator).} There exists a unique self-adjoint operator $\hat{p}$ on $\mathcal{H}$ such that $T_a = \exp(i\hat{p} \cdot a / \hbar)$.
  \item \textbf{F3 (Identification with the momentum operator).} In the position representation of the U2 Hilbert space, $\hat{p} = -i\hbar\nabla$, with eigenfunctions $\langle x | k \rangle = (2\pi\hbar)^{-d/2} e^{ikx/\hbar}$, and the unique unitary intertwiner of position and momentum representations is the standard $L^2$ Fourier transform.
\end{itemize}

\subsection{The four falsifiers}

\begin{itemize}
  \item \textbf{Fal-1 (Non-unitary translations).} $T_a$ acts on $\mathcal{H}$ but does not preserve the U2 inner product.
  \item \textbf{Fal-2 (Multiple inequivalent generators).} The translation group admits more than one self-adjoint generator (Stone's theorem fails on $\mathcal{H}$).
  \item \textbf{Fal-3 (Non-self-adjoint generator).} A generator exists but is not self-adjoint, precluding interpretation as a physical observable.
  \item \textbf{Fal-4 (Non-linear translation action).} The translation action on participation-measure components is non-linear, excluding representation-theoretic treatment.
\end{itemize}

\subsection{Methodological framing: U3 forbidden}

The U5 arc adopts the discipline that \textbf{U3 may not appear in the derivation in any role}: not as a derivation input, not as an existence assumption, not as scope conditioning. This is \emph{not} Framing~1 (which appeared in U4 to handle a sibling-\textsc{candidate} sibling-conditionality structure). U5 is structurally upstream of U3: the spatial translation generator must be installed before the time-translation generator can be identified as its temporal analog (which is U3's content). To derive U5 conditional on U3 would invert the dependency and re-introduce the methodological flaw of the 2026-04-24 tightening-program memo.

The forbidden-input discipline is enforced by an explicit circularity audit (\S\ref{sec:forbidden}).

\section{Structural Inputs}\label{sec:inputs}

\subsection{Primitives}

\begin{itemize}
  \item \textbf{Primitive 03 (Participation):} graph homogeneity. The participation relation does not privilege any vertex of the participation graph. This supplies translation symmetry as a kinematic feature of the graph.
  \item \textbf{Primitive 06 (ED Gradient):} $\nabla \rho$ supplies a continuous spatial axis. Translations along this axis are well-defined and admissible. No preferred origin.
  \item \textbf{Primitive 09 (Tension Polarity):} $U(1)$-valued polarity phase $\pi(K, x)$. The gradient $\nabla \pi$ identifies the phase-propagation rate $k$, which is the wavenumber Fourier-conjugate to position.
  \item \textbf{Primitive 04 (Bandwidth):} supplies the four-band decomposition; the adjacency band $b^{\mathrm{adj}}$ is the band on which the F1--F3 partition acts.
\end{itemize}

\subsection{Inherited \textsc{forced} upstream items}

The U5 arc has access to the following theorem-grade upstream items established in the prior arc cycle:

\begin{itemize}
  \item \textbf{U2-Discrete (Theorem 11).} Sesquilinear inner product on the discrete participation-measure Hilbert space.
  \item \textbf{U2-Continuum (Theorem 12).} Continuum lift of the U2 inner product, with explicit conformal gauge structure.
  \item \textbf{Theorem 10 (Born rule, Gleason--Busch).} Bandwidth-density identifications $b_x \propto |\Psi(x)|^2$ and $b_p \propto |\widetilde{\psi}(p)|^2$ inherit from Theorem~10 applied per band.
\end{itemize}

\subsection{Participation-graph kinematics}

The participation graph $G = (V, E)$ supplies vertices (Primitive~01), edges (Primitive~03), and bandwidth on the edges (Primitive~04). The translation group $\{T_a : a \in \mathbb{R}^d\}$ acts on participation-measure components by $T_a P_K(x) = P_K(x - a)$. This action is well-defined on $\mathcal{H}$ at the kinematic level --- no dynamical evolution required.

\subsection{Mathematical infrastructure}

\begin{itemize}
  \item Stone's theorem on one-parameter unitary groups (Stone 1932)~\cite{stone}.
  \item Standard $L^2$ Fourier transform on $L^2(\mathbb{R}^d)$, with Parseval's identity and the inversion formula~\cite{reedsimon}.
  \item Standard functional calculus on self-adjoint operators.
\end{itemize}

\subsection{Forbidden inputs (circularity audit)}\label{sec:forbidden}

The following inputs are explicitly forbidden as derivation inputs to U5:

\begin{itemize}
  \item \textbf{U3} (the participation-measure evolution equation $i\hbar \partial_t P_K = \hat{H} P_K + \cdots$). Forbidden in any role.
  \item \textbf{The Schr\"odinger equation in textbook form.} Forbidden as a derivation input.
  \item \textbf{Any time-evolution operator or Hamiltonian.} Specifically, $\hat{H}$ does not appear in the U5 derivation.
  \item \textbf{U4} (the specific Hamiltonian form). Forbidden, since U4 is structurally downstream of U5.
\end{itemize}

Each subsequent derivation (\S\S\ref{sec:F1}--\ref{sec:F3}) operates strictly within the primitive-input scope of the present arc and references the forbidden-input list explicitly to confirm acyclicity. This audit follows the canonical templates of U1 Memo~01 \S4 and U4 Memo~01 \S2.

\subsection{Scope conditions}

\begin{itemize}
  \item \textbf{Non-relativistic single-particle scope.} Inherited from Phase-1 Step~2's framing. The U5 derivation does not assume relativistic kinematics; the translation group is the abelian $\mathbb{R}^d$ group, not a Galilean or Lorentzian boost group (those enter U4).
  \item \textbf{Spatial-translation scope.} The U5 arc treats spatial translations only. Time translations (which would invoke Primitive~13 in addition to Primitives~03, 06) are U3's content.
\end{itemize}

\section{Sub-Feature Decomposition}\label{sec:decomposition}

\begin{center}
\begin{tabular}{p{2.5cm}p{1.8cm}p{6.5cm}p{3.5cm}}
\hline
Sub-feature & Status & Structural inputs & Falsifiers \\
\hline
F1 (translation representation) & \textsc{forced} & Primitives 03 + 06 + U2 inner product & Fal-1, Fal-4 \\
F2 (uniqueness of generator) & \textsc{forced} & F1 + Stone's theorem & Fal-2, Fal-3 \\
F3 (momentum identification) & \textsc{forced} & F1 + F2 + U1 polarity-phase + standard $L^2$ Fourier theory & (consequence of F1--F2; structural specificity via dismissal of fractional-Fourier / wavelet / Mellin) \\
\hline
\end{tabular}
\end{center}

All three sub-features are forced-via-derivation. None is automatic; none is load-bearing in the sense of requiring a fresh primitive-level commitment. The closure of F3 requires explicit dismissal of three classes of alternative conjugacies (\S\ref{sec:alternatives}).

\section{Falsifier Inventory}\label{sec:falsifiers}

Each falsifier targets a specific physical alternative and a specific failure mode of the U5 commitment.

\begin{itemize}
  \item \textbf{Fal-1 (Non-unitary translations).} Spatial translations exist as set-theoretic operators on $\mathcal{H}$ but fail to preserve the U2 inner product. Physical interpretation: ``translating'' a state would not preserve probability amplitudes, contradicting the homogeneity of physical space. Resolved in \S\ref{sec:F1} by direct verification that $T_a$ preserves the U2 inner product.
  \item \textbf{Fal-2 (Multiple inequivalent generators).} The translation group admits two or more inequivalent self-adjoint generators on $\mathcal{H}$, leaving the momentum operator structurally ambiguous. Physical interpretation: there would be no unique momentum observable, undermining its operational meaning. Resolved in \S\ref{sec:F2} by Stone's theorem.
  \item \textbf{Fal-3 (Non-self-adjoint generator).} A generator exists but is not self-adjoint. Physical interpretation: momentum would not be a physical observable; eigenvalues need not be real. Resolved in \S\ref{sec:F2}: Stone's theorem produces a self-adjoint generator from a strongly continuous unitary group.
  \item \textbf{Fal-4 (Non-linear translation action).} Translations do not act linearly on participation-measure components, precluding representation-theoretic treatment. Physical interpretation: superposition would not be preserved under translation. Resolved in \S\ref{sec:F1} by linearity of the translation action $T_a P_K(x) = P_K(x - a)$.
\end{itemize}

\section{Derivation of F1 (Translation Representation)}\label{sec:F1}

\subsection{The claim}

The spatial translation group $\{T_a : a \in \mathbb{R}^d\}$ is a well-defined, strongly continuous, abelian $d$-parameter unitary group on the U2 Hilbert space $\mathcal{H}$.

\subsection{Translation symmetry as a kinematic feature}

Define the translation operator $T_a$ acting componentwise on participation-measure components:
\begin{equation}
(T_a P)_K(x) := P_K(x - a), \qquad a \in \mathbb{R}^d. \label{eq:Ta-def}
\end{equation}
Three primitive-level observations:

\begin{itemize}
  \item[\textbf{(i)}] Primitive~03's graph-homogeneity guarantees that no vertex is privileged; the participation relation is invariant under graph automorphisms induced by spatial translations.
  \item[\textbf{(ii)}] Primitive~06's gradient axis $\nabla\rho$ is continuous and unbounded; translations along this axis are well-defined for all $a \in \mathbb{R}^d$.
  \item[\textbf{(iii)}] No primitive-level structure singles out a privileged origin or a privileged translation magnitude.
\end{itemize}

These three observations together establish translation symmetry as a \emph{kinematic} property of the participation graph --- a symmetry of the static structure, not of any dynamical evolution. \textbf{No reference to U3, to time evolution, or to a Hamiltonian appears in this argument.}

\subsection{Linearity}

Translation acts linearly:
\[
T_a (\alpha P + \beta Q)_K(x) = \alpha P_K(x - a) + \beta Q_K(x - a) = \alpha (T_a P)_K(x) + \beta (T_a Q)_K(x)
\]
for all $\alpha, \beta \in \mathbb{C}$ and all $P, Q \in \mathcal{H}$. \textbf{Falsifier Fal-4 dispatched.}

\subsection{Unitarity}

For $P, Q \in \mathcal{H}$ and $a \in \mathbb{R}^d$:
\[
\langle T_a P \mid T_a Q \rangle_{\mathcal{H}} = \sum_K \int_{\mathbb{R}^d} P_K^*(x - a) \, Q_K(x - a) \, d\mu(x).
\]
Substituting $x' = x - a$, with $d\mu(x') = d\mu(x)$ since the U2 measure is translation-invariant (Primitive~06 supplies a translation-invariant Lebesgue measure on the spatial axis; the U2-Continuum gauge acts as a conformal rescaling that does not depend on $x$ in the flat baseline):
\[
\langle T_a P \mid T_a Q \rangle_{\mathcal{H}} = \sum_K \int_{\mathbb{R}^d} P_K^*(x') \, Q_K(x') \, d\mu(x') = \langle P \mid Q \rangle_{\mathcal{H}}.
\]
The translation operators preserve the U2 inner product. \textbf{Falsifier Fal-1 dispatched.}

\subsection{Group structure}

Direct calculation: $T_a T_b = T_{a+b}$, $T_0 = I$, $T_{-a} = T_a^{-1}$. The translation group is abelian and isomorphic to $(\mathbb{R}^d, +)$.

\subsection{Strong continuity}

For any $P \in \mathcal{H}$ with components $P_K$ in $L^2(d\mu)$, the map $a \mapsto T_a P$ is strongly continuous because $L^2$-translation is strongly continuous. (This is a standard property of $L^2$, inherited by the direct sum / direct integral structure of the U2 Hilbert space across the channel index $K$.)

\begin{theorem}[F1 --- Translation Representation]
Let $\mathcal{H}$ denote the U2 Hilbert space of participation measures established by the U2-Discrete and U2-Continuum theorems. Then the spatial translation operators $T_a P_K(x) := P_K(x - a)$ for $a \in \mathbb{R}^d$ form a strongly continuous abelian $d$-parameter unitary group on $\mathcal{H}$, derived from Primitives~03 and 06 plus the U2 inner-product structure, with no input from U3 or any time-evolution operator.
\end{theorem}

\section{Derivation of F2 (Uniqueness of the Generator)}\label{sec:F2}

\subsection{Stone's theorem}

\begin{theorem}[Stone 1932~\cite{stone}]
Every strongly continuous one-parameter unitary group $\{U(t) : t \in \mathbb{R}\}$ on a Hilbert space has a unique self-adjoint generator $A$ such that $U(t) = \exp(itA)$.
\end{theorem}

\subsection{Application to F1's translation group}

By Theorem~F1, $\{T_a : a \in \mathbb{R}^d\}$ is a strongly continuous unitary group on $\mathcal{H}$. Apply Stone's theorem to each spatial direction $\hat{e}_i$ with $a = a_i \hat{e}_i$: there exists a unique self-adjoint generator $\hat{p}_i$ on $\mathcal{H}$ such that
\[
T_{a_i \hat{e}_i} = \exp\!\left(\frac{i \hat{p}_i a_i}{\hbar}\right).
\]
The $d$ generators $\hat{p}_i$ commute pairwise (by the abelian structure of the translation group, $T_a T_b = T_{a+b} = T_b T_a$), so the joint exponentiation is well-defined:
\begin{equation}
T_a = \exp\!\left(\frac{i \hat{p} \cdot a}{\hbar}\right), \qquad \hat{p} := (\hat{p}_1, \ldots, \hat{p}_d). \label{eq:Ta-stone}
\end{equation}
The generator $\hat{p}$ is unique. \textbf{Falsifier Fal-2 dispatched.}

\subsection{Self-adjointness}

Stone's theorem produces a self-adjoint generator from a strongly continuous unitary group. The self-adjointness is automatic from the theorem and does not require separate verification. \textbf{Falsifier Fal-3 dispatched.}

\subsection{Why this does not require U3}

Stone's theorem is a statement about strongly continuous one-parameter unitary groups on Hilbert spaces. The translation group $\{T_a\}$ is one such group on $\mathcal{H}$. The hypotheses of Stone's theorem are entirely kinematic: the existence of $\mathcal{H}$ (supplied by U2), the unitarity of $T_a$ (Theorem~F1), and the strong continuity of $T_a$ (Theorem~F1). At no point does the theorem require the existence of a \emph{time}-evolution operator, a Hamiltonian, or any dynamical input. The translation generator $\hat{p}$ is identified as a \emph{spatial} self-adjoint operator on $\mathcal{H}$ at fixed time, with no reference to the participation-measure evolution equation.

\begin{theorem}[F2 --- Uniqueness of the Translation Generator]
Under the hypotheses of Theorem~F1, there exists a unique self-adjoint operator $\hat{p}$ on $\mathcal{H}$ such that $T_a = \exp(i\hat{p} \cdot a / \hbar)$ for all $a \in \mathbb{R}^d$. The identification is via Stone's theorem applied to the strongly continuous unitary group $\{T_a\}$, and is U3-independent.
\end{theorem}

\section{Derivation of F3 (Momentum Operator Identification)}\label{sec:F3}

\subsection{The position-representation form}

In the position representation of $\mathcal{H}$, the translation operator acts on a wavefunction $\psi(x)$ by $T_a \psi(x) = \psi(x - a)$. Taylor-expansion in $a$ for analytic $\psi$:
\begin{equation}
T_a \psi(x) = \psi(x - a) = \sum_{n=0}^{\infty} \frac{(-a \cdot \nabla)^n}{n!} \psi(x) = \exp(-a \cdot \nabla) \psi(x). \label{eq:Ta-taylor}
\end{equation}
Comparing with~\eqref{eq:Ta-stone}:
\[
\exp\!\left(\frac{i \hat{p} \cdot a}{\hbar}\right) = \exp(-a \cdot \nabla)
\]
which is satisfied by
\begin{equation}
\boxed{\hat{p} = -i\hbar\nabla.} \label{eq:p-position}
\end{equation}

\subsection{Plane-wave eigenfunctions}

The eigenvalue equation $\hat{p} \, |k\rangle = k |k\rangle$ in the position representation reads $-i\hbar\nabla \langle x | k \rangle = k \langle x | k \rangle$, with solution
\begin{equation}
\langle x | k \rangle = (2\pi\hbar)^{-d/2} \, e^{i k \cdot x / \hbar}. \label{eq:planewave}
\end{equation}
The normalization $(2\pi\hbar)^{-d/2}$ is fixed by the orthonormality $\langle k | k' \rangle = \delta(k - k')$ via the standard Fourier convention.

\subsection{The thin-limit / continuum identification}

In the discrete regime, the channel index $K$ runs over a discrete set; in the continuum / thin limit, $K$ is replaced by a continuous wavenumber $k \in \mathbb{R}^d$. The identification is:
\begin{equation}
K \xleftrightarrow{\text{thin limit}} k. \label{eq:Kk}
\end{equation}
The U1 polarity-phase factor $e^{i\pi_K(x)}$ in the participation-measure construction $P_K = \sqrt{b_K} \, e^{i\pi_K}$ is what makes this identification natural: in the continuum limit, $\pi_K(x) \to k \cdot x / \hbar$, so the polarity phase is precisely the plane-wave phase. The U1-supplied phase structure is the structural bridge between the channel index and the momentum eigenvalue.

\subsection{The role of the U2 inner product}

The U2 inner product $\langle P | Q \rangle = \sum_K \int P_K^* Q_K \, d\mu$ supplies the Hilbert-space structure on which:

\begin{itemize}
  \item[\textbf{(i)}] The Fourier transform $\widetilde{\psi}(k) = (2\pi\hbar)^{-d/2} \int e^{-ikx/\hbar} \psi(x) \, dx$ is an isometry: $\langle \psi | \phi \rangle_x = \langle \widetilde{\psi} | \widetilde{\phi} \rangle_k$ (Parseval).
  \item[\textbf{(ii)}] The position eigenstates $\{|x\rangle\}$ and momentum eigenstates $\{|k\rangle\}$ are both complete orthonormal bases.
  \item[\textbf{(iii)}] The standard Fourier transform is the unique unitary intertwiner between the two bases.
\end{itemize}

\subsection{Dismissal of alternative conjugacies}\label{sec:alternatives}

Three classes of alternative integral transforms might in principle have served as the conjugacy between position and momentum representations. Each is dismissed:

\begin{itemize}
  \item \textbf{Fractional-Fourier transforms.} Diagonalize rotated phase-space operators (specifically, the operator $\cos\theta \cdot \hat{x} + \sin\theta \cdot \hat{p}$ for fractional angle $\theta$). Not the translation generator $\hat{p}$. Fail to provide the \emph{translation}-conjugate basis. Inconsistent with the Stone's-theorem identification.
  \item \textbf{Wavelet transforms.} Multi-scale resolutions; not Stone's-theorem intertwiners of any one-parameter unitary group. Provide a different mathematical structure (overcomplete frames, time-frequency tilings) that is not the translation-conjugate basis.
  \item \textbf{Mellin transforms.} Intertwine the dilation group $\{x \to e^t x\}$, not the translation group. Dilations are not a primitive-level kinematic symmetry of the participation graph (the participation graph has no preferred scale; dilation symmetry is not supplied by any primitive). Mellin transforms therefore do not arise structurally.
\end{itemize}

The standard $L^2$ Fourier transform is the \textbf{unique} unitary intertwiner of position and translation-conjugate momentum representations. It is forced by (a) the translation symmetry of the participation graph, (b) the Stone's-theorem identification of the unique self-adjoint generator, and (c) the standard $L^2$ Fourier theory.

\begin{theorem}[F3 --- Momentum Operator Identification]
Let $\hat{p}$ be the unique self-adjoint translation generator on the U2 Hilbert space $\mathcal{H}$ supplied by Theorem~F2. Then in the position representation of $\mathcal{H}$, $\hat{p} = -i\hbar\nabla$. The eigenfunctions of $\hat{p}$ are plane waves $\langle x | k \rangle = (2\pi\hbar)^{-d/2} \, e^{ikx/\hbar}$. The unique unitary intertwiner of position and momentum representations is the standard $L^2$ Fourier transform; fractional-Fourier, wavelet, and Mellin alternatives are dismissed by structural inconsistency with the translation-driven conjugacy.
\end{theorem}

\section{Falsifier Resolution Table}\label{sec:falsifier-table}

\begin{center}
\begin{tabular}{p{2.0cm}p{1.0cm}p{8.5cm}p{1.8cm}}
\hline
Falsifier & Target & Dispatch argument & Section \\
\hline
Fal-1 & F1 & Direct calculation: $\langle T_a P \mid T_a Q \rangle = \langle P \mid Q \rangle$ via translation-invariance of the U2 measure & \S\ref{sec:F1}.4 \\
Fal-2 & F2 & Stone's theorem on the strongly continuous unitary group $\{T_a\}$ produces a unique self-adjoint generator & \S\ref{sec:F2}.2 \\
Fal-3 & F2 & Stone's theorem produces a self-adjoint generator from a strongly continuous unitary group & \S\ref{sec:F2}.3 \\
Fal-4 & F1 & Direct linearity of $T_a$ on participation-measure components & \S\ref{sec:F1}.3 \\
\hline
\end{tabular}
\end{center}

\textbf{All four falsifiers dispatched within the stated scope conditions and forbidden-input discipline.}

\section{Conditionality Ledger}\label{sec:conditionality}

The U5 verdict is \textbf{\textsc{forced}} with the following inherited conditionalities and zero new ones:

\begin{enumerate}
  \item \textbf{Non-relativistic single-particle scope.} Inherited from Phase-1 Step~2; conditions the use of the abelian $(\mathbb{R}^d, +)$ translation group rather than a Galilean or Lorentzian boost group.
  \item \textbf{U2 structure (inner product, both regimes).} Inherited from the U2-Discrete and U2-Continuum theorems; supplies the Hilbert space on which Stone's theorem applies.
  \item \textbf{U1 participation-measure construction.} Inherited from the U1 theorem (\textsc{forced}); supplies the polarity-phase structure that natural-identifies the channel index with the momentum eigenvalue in the continuum limit (\S\ref{sec:F3}.3).
  \item \textbf{Theorem~10 (Born) for downstream use only.} Used in the bandwidth-density identifications $b_x \propto |\Psi|^2$, $b_p \propto |\widetilde{\psi}|^2$ in the equivalence to the adjacency-band-partition framing; not used in the F1--F3 derivations themselves.
\end{enumerate}

\textbf{No dependence on U3.} The U5 derivation invokes no time-evolution operator, no Hamiltonian, no Schr\"odinger equation. The translation symmetry is kinematic (Primitives~03 + 06), the unitary representation is established at fixed time on the U2 Hilbert space (Theorem~F1), Stone's theorem identifies the generator without dynamical input (Theorem~F2), and the position-representation form (Theorem~F3) follows from comparison of operator exponentials. \textbf{U5 is structurally upstream of U3 and remains so under the present derivation.}

\textbf{No new structural \textsc{candidate}s introduced.} All four conditionalities are inherited from prior structural commitments.

\section{Final U5 Theorem}\label{sec:final-theorem}

\begin{theorem}[U5; Translation Symmetry and the Momentum Operator]
Let the ED primitive stack supply: graph homogeneity (Primitive~03), bandwidth (Primitive~04), spatial gradient (Primitive~06), and tension polarity (Primitive~09). Let the participation-measure carrier (U1, \textsc{forced}) and the U2 inner product on the participation-measure Hilbert space (U2-Discrete + U2-Continuum, both \textsc{forced}) supply the Hilbert space $\mathcal{H}$. Then:

\begin{enumerate}
  \item \textbf{Translation representation.} The spatial translation operators $T_a P_K(x) := P_K(x - a)$ for $a \in \mathbb{R}^d$ form a strongly continuous abelian $d$-parameter unitary group on $\mathcal{H}$.

  \item \textbf{Unique self-adjoint generator.} By Stone's theorem, there exists a unique self-adjoint operator $\hat{p}$ on $\mathcal{H}$ such that
  \[
  T_a = \exp\!\left(\frac{i \hat{p} \cdot a}{\hbar}\right), \qquad a \in \mathbb{R}^d.
  \]

  \item \textbf{Momentum operator identification.} In the position representation of $\mathcal{H}$, $\hat{p} = -i\hbar\nabla$. The eigenfunctions of $\hat{p}$ are plane waves $\langle x | k \rangle = (2\pi\hbar)^{-d/2} e^{ikx/\hbar}$, with the standard $L^2$ Fourier transform as the unique unitary intertwiner of position and momentum representations.
\end{enumerate}

The derivation invokes no U3 input, no Hamiltonian, no time-evolution operator, and no Schr\"odinger equation. All four falsifiers (non-unitary translations, multiple inequivalent generators, non-self-adjoint generator, non-linear translation action) are dispatched.
\end{theorem}

\textbf{Status:} \textsc{forced}, unconditional in both discrete and continuum regimes; gauge-compatible with the U2-Continuum conformal redundancy (the gauge acts uniformly on both partition components and trivially on the translation representation); U3-independent.

\section{Downstream Consequences}\label{sec:downstream}

\subsection{How U5 feeds into U4 (momentum identification $\to$ kinetic term)}

U4's Hamiltonian form $\hat{H} = \hbar^2 |\hat{p}|^2 / (2m) + V(\hat{x})$ requires the prior identification of the momentum operator $\hat{p}$. U5 supplies $\hat{p}$ as a self-adjoint operator on the U2 Hilbert space with plane-wave eigenfunctions. U4's Galilean Lie algebra integration acts on this $\hat{p}$:
\[
[\hat{H}, \hat{K}_i] = -i\hbar \hat{p}_i
\]
with $\hat{K}_i = m\hat{x}_i - t\hat{p}_i$. The integration of this commutator condition (yielding $\partial T / \partial \hat{p}_i = \hat{p}_i / m$ and thus $T = |\hat{p}|^2 / (2m)$) presupposes that $\hat{p}_i$ is a well-defined self-adjoint operator. \textbf{U5 supplies that operator.} Without U5, the U4 Galilean integration is not statable.

\subsection{How U5 feeds into U1 (complex structure compatibility)}

The U1 participation-measure construction $P_K = \sqrt{b_K} \cdot e^{i\pi_K}$ supplies the polarity-phase factor $e^{i\pi_K}$ whose gradient identifies the wavenumber $k = \nabla \pi$. U5 closes the loop: the wavenumber $k$ is the eigenvalue of the translation generator $\hat{p}$, and the complex-valued structure of the participation measure is precisely what makes the plane-wave eigenfunctions $e^{ikx/\hbar}$ available as a basis. The U1 complex structure and the U5 momentum operator are mutually compatible --- neither was designed assuming the other, but both close consistently within the same Hilbert-space framework.

\subsection{Updated structural-foundations ledger}

Following the 2026-04-26 cycle, the inventory of structural-foundations theorems stands at fifteen forced items, with U5 occupying position \#13:

\begin{center}
\begin{tabular}{cll}
\hline
\# & Theorem & Status \\
\hline
10 & Born rule (Gleason--Busch) & \textsc{forced} \\
11 & U2-Discrete (sesquilinear inner product) & \textsc{forced} \\
12 & U2-Continuum (with conformal gauge) & \textsc{forced} \\
\textbf{13} & \textbf{U5 (translation generator + momentum operator; this paper)} & \textbf{\textsc{forced}} \\
14 & U1 (participation-measure construction) & \textsc{forced} \\
15 & U4 (specific Hamiltonian form; form-\textsc{forced}--value-\textsc{inherited}) & \textsc{forced} cond.\ on U3 (Framing 1) \\
\hline
\end{tabular}
\end{center}

\subsection{Forced theorem count and active CANDIDATE inventory}

The U5 paper formalizes a result already promoted in the 2026-04-26 cycle. The active upstream-\textsc{candidate} inventory therefore reads, post-cycle:
\[
\{U3\} + \text{continuum gauge convention.}
\]
One upstream \textsc{candidate} remains. Three of the four foundational QM postulates (Born rule, Bell--Tsirelson bound, Heisenberg uncertainty) are \textsc{forced}-unconditional. Schr\"odinger evolution is \textsc{forced} conditional on U3 alone.

\section{Discussion and Outlook}\label{sec:discussion}

\subsection{Why U5 is the kinematic backbone of the structural-foundations program}

The structural-foundations cycle decomposes into three layers:

\begin{itemize}
  \item \textbf{Layer 1 (carrier):} U1 --- the participation-measure construction.
  \item \textbf{Layer 2 (Hilbert-space arena):} U2-Discrete + U2-Continuum + Theorem~10 + \textbf{U5}.
  \item \textbf{Layer 3 (dynamical):} U3 + U4.
\end{itemize}

U5 occupies the Hilbert-space layer alongside U2 and Theorem~10. Within that layer, U5 is what installs the \emph{kinematic} generators: spatial translations have a unique self-adjoint generator (the momentum operator), with plane-wave eigenfunctions and a well-defined Fourier-conjugate position basis. Without U5, the Hilbert space exists but carries no identified momentum observable; with U5, the kinematic skeleton of standard quantum mechanics is in place at fixed time.

The ``kinematic backbone'' framing is precise: U5 supplies everything needed to do single-time quantum kinematics --- observables, eigenbases, conjugate-pair structure, Fourier conjugacy --- without invoking time evolution. Time evolution is U3's content; the Hamiltonian's specific form is U4's content. U5 is the structural boundary between kinematics and dynamics.

\subsection{How U5 + U2 + U4 jointly force the Schr\"odinger form}

The dynamical content of non-relativistic single-particle quantum mechanics is the Schr\"odinger equation $i\hbar \partial_t \psi = \hat{H} \psi$ with $\hat{H} = \hbar^2|\hat{p}|^2/(2m) + V(\hat{x})$. Three structural commitments suffice to force it:

\begin{itemize}
  \item \textbf{U2} supplies the Hilbert space.
  \item \textbf{U5} supplies the momentum operator $\hat{p}$ on that Hilbert space.
  \item \textbf{U4} supplies the Hamiltonian form $\hat{H} = \hbar^2|\hat{p}|^2/(2m) + V(\hat{x})$ via Galilean Lie algebra integration.
\end{itemize}

The remaining commitment --- U3, the existence and linear-first-order structure of the time-evolution generator --- is what binds these three pieces into a dynamical equation. With U3 closed, the Schr\"odinger equation is structurally complete: U2 supplies the arena, U5 supplies the kinematic generator, U4 supplies the dynamical generator's form, U3 supplies the existence and linearity of dynamical evolution.

The U5 derivation is what makes the U4 closure acyclic: U4's Galilean integration acts on the U5 momentum operator, and U5 was derived without invoking U4 or U3. Promoting U3 in the next arc would close the entire QM-emergence Phase-1 program.

\subsection{Sharpening over earlier work}

The earlier 2026-04-24 tightening-program memo \texttt{arcs/arc-foundations/u5\_adjacency\_partition\_derivation.md} derived the Fourier conjugacy via U3 --- invoking the participation-measure evolution equation as a derivation input. The present derivation removes that dependency: translation symmetry is established as a kinematic property of the participation graph (Primitives~03 + 06), Stone's theorem applies on the U2 Hilbert space without time-evolution input, and the Fourier conjugacy follows from the position-representation form of the translation generator together with the standard $L^2$ Fourier theory.

The methodological consequence is decoupling: U5's verdict is independent of whatever conditionality U3 carries. Promoting U3 in a future arc does not retroactively change the U5 result, and U5's promotion does not feed back into U3's derivation chain.

\subsection{Next steps}

\begin{enumerate}
  \item \textbf{Open the U3 arc.} The natural and final foundational target. U3 inherits all six \textsc{forced} upstream items (U1, U2-Discrete, U2-Continuum, U5, Theorem~10, U4) and applies Stone's theorem to time-translation symmetry (Primitive~13) --- paralleling F1 of the present paper applied to spatial translations. Headline result if it closes: all four foundational QM postulates structurally derived from ED primitives.
  \item \textbf{Synthesis paper revision.} With U5 now formalized as a publication paper, the QM-emergence Phase-1 synthesis paper \S\S3.3/3.4/3.5/4 can be revised to reflect the four-paper publication series (Born\_Gleason, U2\_Inner\_Product, U1\_Participation\_Measure, U4\_Hamiltonian\_Form) plus this paper. Recommend deferring the synthesis revision until U3 closes (single editorial pass).
  \item \textbf{Public-facing U5 explainer.} A natural fifth in the public-explainer ladder, sitting between the U2 explainer (Hilbert space) and the U4 explainer (Hamiltonian form). The arc closure memo \S6 supplies a draft.
\end{enumerate}

\begin{thebibliography}{99}

\bibitem{phase1}
A.~Proxmire, \emph{Phase-1: QM Emergence Structural Completion}, \texttt{papers/QM\_Emergence\_Structural\_Completion/}, 2026.

\bibitem{borngleason}
A.~Proxmire, \emph{The Born Rule as a Forced Theorem of Event Density: A Gleason--Busch Reconstruction from First Principles}, \texttt{papers/Born\_Gleason/}, 2026.

\bibitem{u2}
A.~Proxmire, \emph{The Inner Product as Forced Structure in Event Density: Discrete Derivation, Continuum Lift, and Gauge-Invariant Completion}, \texttt{papers/U2\_Inner\_Product/}, 2026.

\bibitem{u1}
A.~Proxmire, \emph{The U1 Theorem: Deriving the Participation Measure from Primitive Structure}, \texttt{papers/U1\_Participation\_Measure/}, 2026.

\bibitem{u4}
A.~Proxmire, \emph{U4: The Forced Structure of the Non-Relativistic Hamiltonian}, \texttt{papers/U4\_Hamiltonian\_Form/}, 2026.

\bibitem{arcm}
A.~Proxmire, \emph{Arc M: Chain-Mass Structure in Event Density}, \texttt{papers/Arc\_M/}, 2026.

\bibitem{u5arc}
A.~Proxmire, \emph{U5 arc memos (decomposition, F1+F2+F4 derivations, F3+F5 verdict, closure)}, \texttt{arcs/U5/}, 2026.

\bibitem{stone}
M.~H.~Stone, ``On one-parameter unitary groups in Hilbert space,'' \emph{Annals of Mathematics} \textbf{33}, 643--648 (1932).

\bibitem{bargmann}
V.~Bargmann, ``On unitary ray representations of continuous groups,'' \emph{Annals of Mathematics} \textbf{59}, 1--46 (1954).

\bibitem{reedsimon}
M.~Reed and B.~Simon, \emph{Methods of Modern Mathematical Physics, Volume I: Functional Analysis}, Academic Press, 1980.

\bibitem{prugovecki}
E.~Prugove\v{c}ki, \emph{Quantum Mechanics in Hilbert Space}, 2nd ed., Academic Press, 1981.

\end{thebibliography}

\vspace{1em}
\noindent\rule{\linewidth}{0.4pt}
\vspace{0.2em}

\noindent\emph{Manuscript prepared April 2026. Source materials: \texttt{arcs/U5/} four-memo arc and closure memo, formalized into a single publishable document. The U5 arc is the third of the structural-foundations cycle (chronologically the fourth, drafted as the fifth publication paper after Born\_Gleason, U2\_Inner\_Product, U1\_Participation\_Measure, and U4\_Hamiltonian\_Form).}

\noindent\emph{Draft complete. Pausing for revision.}

\end{document}
